\section*{Hoofdstuk \ref{ch:g11:absorbance} --- Kleur en absorbantie}

\begin{solution}{exo:g11:absorbance:1}
Als ze oranje opneemt, ziet ze er blauw uit, de kleur tegenover oranje op
de cirkel. Als ze violet opneemt, ziet ze er geel uit.
\end{solution}

\begin{solution}{exo:g11:absorbance:2}
Een groene oplossing neemt vooral rood licht op, de kleur tegenover groen
op de cirkel: van ongeveer \qty{620}{nm} tot \qty{750}{nm}.
\end{solution}

\begin{solution}{exo:g11:absorbance:3}
Blauwe kleurstof: $\lambda_{\max} = \qty{629}{nm}$, $A = 0.82$. Gele
kleurstof: $\lambda_{\max} = \qty{427}{nm}$, $A = 0.795$, ongeveer 0.80.
\end{solution}

\begin{solution}{exo:g11:absorbance:4}
$A$ is evenredig met de concentratie: drie keer zoveel geeft
$A = 0.90$; half zoveel geeft $A = 0.15$.
\end{solution}

\begin{solution}{exo:g11:absorbance:5}
De blanco is een cuvet gevuld met het zuivere \omterm{def:g6:solutions-and-solubility:solute}{oplosmiddel}. Door er
$A = 0$ mee in te stellen, haal je weg wat de cuvet, het \omterm{def:g6:solutions-and-solubility:solute}{oplosmiddel} en
het instrument opnemen, zodat de metingen daarna alleen van de \omterm{def:g6:solutions-and-solubility:solute}{opgeloste
stof} komen.
\end{solution}

\begin{solution}{exo:g11:absorbance:6}
$C = 0.41 / 0.164 = \qty{2.5}{mg/L}$.
\end{solution}

\begin{solution}{exo:g11:absorbance:7}
Bij \qty{629}{nm} is de \omterm{def:g11:absorbance:absorbance}{absorbantie} het grootst, dus is de meting het
gevoeligst (kleine concentraties geven nog leesbare \omterm{def:g11:absorbance:absorbance}{absorbanties}); en
bovenop de band is de kromme vlak, zodat een kleine fout in de golflengte
de meting nauwelijks verandert. Bij \qty{550}{nm} is de \omterm{def:g11:absorbance:absorbance}{absorbantie}
klein en verandert ze snel met de golflengte.
\end{solution}

\begin{solution}{exo:g11:absorbance:8}
Verdund: $C = 0.48 / 0.164 = \qty{2.9}{mg/L}$; de drank:
$5 \times 2.9 = \qty{15}{mg/L}$ (preciezer $14.6$).
\end{solution}

\begin{solution}{exo:g11:absorbance:9}
$a = A / (\ell\, C_m) = 0.795 / (1.00 \times 0.0150) = \qty{53.0}{L/(g.cm)}$;
$\varepsilon = a \times M = 53.0 \times 534.4 = \qty{2.83e4}{L/(mol.cm)}$.
\end{solution}

\begin{solution}{exo:g11:absorbance:10}
Boven ongeveer \qty{520}{nm}. Bij \qty{629}{nm} neemt de gele kleurstof
helemaal niets op, dus komt de hele \omterm{def:g11:absorbance:absorbance}{absorbantie} daar van de blauwe
kleurstof.
\end{solution}

\begin{solution}{exo:g11:absorbance:11}
Ze verdubbelt: $A = \varepsilon \ell c$ is evenredig met $\ell$; het
licht gaat door twee keer zoveel oplossing, dus langs twee keer zoveel
absorberende \omterm{def:g7:atoms-and-molecules:molecule}{moleculen}.
\end{solution}

\begin{solution}{exo:g11:absorbance:12}
$A/C$: $0.164$, $0.162$, $0.115$, $0.0675$. De verhouding is alleen voor
de eerste twee constant: boven ongeveer \qty{10}{mg/L} (\omterm{def:g11:absorbance:absorbance}{absorbanties}
boven ongeveer 1.6) geldt de wet niet meer. Alleen de standaarden van 5
en \qty{10}{mg/L} (en lagere) mag je gebruiken; een onbekende die 2.5
geeft, moet verdund worden, bijvoorbeeld vijf keer, en opnieuw gemeten.
\end{solution}

\begin{solution}{exo:g11:absorbance:13}
Blauwe kleurstof, uit \qty{629}{nm} alleen: $C = 0.574 / 0.164 = \qty{3.5}{mg/L}$.
Gele kleurstof, uit \qty{427}{nm}: $C = 0.53 / 0.0530 = \qty{10}{mg/L}$.
\end{solution}

\begin{solution}{exo:g11:absorbance:14}
Ze vindt $0.368 / 0.164 = \qty{2.24}{mg/L}$. De echte \omterm{def:g11:absorbance:absorbance}{absorbantie} is
$0.368 - 0.040 = 0.328$, dus de echte concentratie is
$0.328 / 0.164 = \qty{2.00}{mg/L}$. Ze zit \qty{12}{\percent} te hoog.
\end{solution}

\begin{solution}{exo:g11:absorbance:15}
$10 \times 60 = \qty{600}{mg}$ per dag, in $600 / 20 = \qty{30}{L}$
drank. Niemand drinkt dertig liter per dag: de drank alleen is geen
realistisch risico, al kan de kleurstof ook uit ander voedsel komen.
\end{solution}

\begin{solution}{pb:g11:absorbance:1}
\textbf{1.} Oranjerood, de kleur tegenover blauw op de cirkel.

\textbf{2.} $\lambda_{\max} = \qty{629}{nm}$.

\textbf{3.} Bij \qty{450}{nm} is het licht blauw: de kleurstof laat het
door, en juist daarom ziet de drank er blauw uit.

\textbf{4.} Bij \qty{629}{nm}.

\textbf{5.} $\qty{50.0}{mg} / \qty{0.5000}{L} = \qty{100}{mg/L}$.

\textbf{6.} \omterm{def:g10:concentration-and-dilution:dilution}{Verdunningsfactor} $100 / 3.0 = 33.3$, dus
$V_0 = 100.0 / 33.3 = \qty{3.00}{mL}$: een maatpipet (of een buret) en
een maatkolf van \qty{100.0}{mL}.

\textbf{7.} $A/C$ = 0.168, 0.163, 0.165, 0.163, 0.165: allemaal dicht bij
0.164, dus $A \approx 0.164\, C$, een lijn door de oorsprong.

\textbf{8.} Een oplossing zonder kleurstof ($C = 0$) is de blanco, die op
$A = 0$ is ingesteld.

\textbf{9.} $a = \qty{0.164}{L/mg}$ per centimeter, dat is
\qty{164}{L/(g.cm)}; $\varepsilon = 164 \times 792.9 =
\qty{1.30e5}{L/(mol.cm)}$.

\textbf{10.} $F = 50.0 / 25.0 = 2$.

\textbf{11.} $C = 0.590 / 0.164 = \qty{3.60}{mg/L}$.

\textbf{12.} $2 \times 3.60 = \qty{7.20}{mg/L}$.

\textbf{13.} Ongeveer $2 \times 0.590 = 1.18$, boven de hoogste
standaard (0.823): de meting zou buiten het geijkte gebied liggen, waar
de wet niet meer hoeft te gelden. Door te verdunnen komt ze tussen de
standaarden.

\textbf{14.} $\qty{7.20}{mg/L} \times \qty{0.500}{L} = \qty{3.60}{mg}$.

\textbf{15.} $n = \num{3.60e-3} / 792.9 = \qty{4.54e-6}{mol}$.

\textbf{16.} $6 \times 30 = \qty{180}{mg}$ per dag.

\textbf{17.} $180 / 3.60 = 50$ flessen.

\textbf{18.} $50 \times 0.500 = \qty{25}{L}$ op één dag: onmogelijk om te
drinken. De kleurstof van deze drank alleen kan een kind niet in de buurt
van de grens brengen, al tellen de kleurstoffen uit verschillende
voedingsmiddelen op.

\textbf{19.} $6 \times 70 = \qty{420}{mg}$, dat is $420 / 3.60 \approx
117$ flessen.

\textbf{20.} Ongeveer 50 flessen van \qty{500}{mL}, \qty{25}{L} drank, op
één dag.
\end{solution}
